% ==========================================================
\section{结论}
\label{section8}
% =========================================================

\subsection{主要研究成果}

本文针对上海市高校甲型流感H3N2的传播特征与防控需求，建立了两群体分层SEIQR传播动力学模型，并完成了参数估计、敏感性分析和防控策略优化．主要研究成果如下：

\textbf{（1）模型构建方面}

本文建立了符合上海高校人口结构和接触特征的传播动力学模型，主要创新点包括：

\begin{itemize}
  \item \textbf{两群体分层结构}：将校园人群划分为学生和教职工两个群体，通过场所加权接触矩阵刻画宿舍、教室、课外活动等场景的异质性接触结构，避免了传统均质混合模型的系统性偏差；
  
  \item \textbf{双谐波季节强迫函数}：采用年周期和半年周期两个分量的叠加，准确捕捉了上海地区冬夏双峰的流行特征，拟合优度$R^2$达到\textbf{0.54}；
  
  \item \textbf{学期调制机制}：引入学期调制系数反映寒暑假期间接触率的骤降，成功捕捉了"开学效应"，为预测秋季开学初的聚集性疫情风险提供了依据；
  
  \item \textbf{SEIQR舱室结构}：引入隔离舱室Q，实现了对校园隔离机制的参数化建模，为量化隔离措施的效果提供了基础．
\end{itemize}

\textbf{（2）参数估计方面}

本文基于多层次数据源完成了参数估计，主要结果包括：

\begin{itemize}
  \item 基础传播系数$\beta_0 =$ \textbf{0.041}（天$^{-1}$），通过WHO FluNet中国H3N2数据拟合得到；
  \item 基本再生数$R_0 =$ \textbf{1.14}（95\% CI: \textbf{1.10—1.18}），与文献报道的季节性H3N2中位值（$R_0=1.25$）基本一致；
  \item 群体免疫阈值$p^* =$ \textbf{12.5\%}，显著高于上海高校当前的疫苗接种率（0.32\%）；
  \item 季节振幅参数$\delta_1 =$ \textbf{0.26}，$\delta_2 =$ \textbf{0.80}，$\delta_2 > \delta_1$表明半年周期分量主导季节波动，冬夏双峰由两分量叠加共同决定．
\end{itemize}

参数估计的相对不确定性约为10\%—20\%，在可接受范围内，模型验证结果显示RMSE = \textbf{0.018}，MAPE = \textbf{46.9\%}，$R^2 =$ \textbf{0.539}，趋势拟合良好（Pearson相关系数$r=0.97$）；MAPE偏高源于非流行期阳性率接近零导致的分母敏感性，需结合RMSE综合判断．

\textbf{（3）敏感性分析方面}

本文采用PRCC方法进行了全局敏感性分析，主要发现包括：

\begin{itemize}
  \item 半年季节振幅$\delta_2$对累计感染规模影响最显著（PRCC$=0.90$，$p<0.001$），表明季节强迫是总感染规模的首要决定因素；$\beta_0$、$\delta_1$和$\gamma$的影响次之（$|$PRCC$|\approx0.18$—$0.19$，$p<0.05$）；
  \item 康复率$\gamma$对峰值感染规模影响最大（PRCC$=-0.33$，$p<0.001$），隔离成功比例$p_{\text{iso}}$次之（PRCC$=0.29$，$p<0.001$），反映了缩短传染期和强化隔离的有效性；
  \item 疫苗接种率$v$在当前水平（0.32\%）下的影响最弱（PRCC$\approx0.01$，$p>0.8$），表明低接种率背景下需大幅提升至临界阈值（12.5\%）才能产生显著效果．
\end{itemize}

敏感性分析为防控资源的优先配置提供了理论依据：应重点关注传播系数的降低（口罩、通风）和接触率的减少（线上授课、社团限流）．

\textbf{（4）防控策略优化方面}

本文构建了含感染总量、防控成本、教学影响的加权多目标优化模型，通过网格搜索给出了分场景最优方案：

\begin{itemize}
  \item \textbf{常态场景}：以疫苗接种推广和教室通风为主，峰值感染率降低约\textbf{25\%—37\%}（冬季至夏季）；
  \item \textbf{散发场景}：全面启动口罩、通风、隔离和社团限流，预期累计感染降低约\textbf{41\%—64\%}；
  \item \textbf{聚集场景}：全面启动所有措施，累计感染可降低约\textbf{63\%—69\%}，峰值感染率降低约\textbf{43\%}．
\end{itemize}

优化方案具有较强的可操作性，可为高校管理者和疾控部门提供定量决策支持．

\subsection{主要创新点}

本文的主要创新点可归纳为以下二点：

\textbf{创新点一：建立了符合上海高校特征的传播动力学模型}

本文首次将场所加权接触矩阵引入校园流感传播模型，通过宿舍、教室、课外活动的权重分解，实现了对校园异质性接触结构的精准刻画．相比传统均质混合模型，本文模型在峰值时间、峰值规模、流行持续时长等关键指标上的预测精度显著提升．

\textbf{创新点二：构建了多目标防控策略优化模型}

本文首次将校园流感防控策略的选择建模为多目标优化问题，在感染控制、经济成本、教学影响三个目标间寻求平衡，给出了分场景、可落地的最优方案．相比经验性决策，优化模型能够提供定量化的策略建议，提高决策的科学性．

\subsection{研究局限与展望}

本文的研究存在以下局限，需要在后续研究中改进：

\textbf{数据层面的局限}：本文使用WHO FluNet全国数据进行参数估计和验证，与校园特定人群存在尺度差异．未来可与上海高校校医院合作，获取校园级别的流感监测数据，提高模型的区域适用性．

\textbf{模型层面的局限}：均质混合假设忽略了宿舍楼层、专业院系等次级结构，可能低估局部聚集爆发的速度．未来可引入网络传播模型或Agent-based模型，刻画个体层面的接触网络和传播过程．

\textbf{干预措施层面的局限}：干预措施的参数化基于文献值和理论推导，未考虑措施实施的动态过程和依从性问题．未来可结合问卷调研和行为监测数据，估计实际的措施依从率，并引入依从性衰减模型．

\textbf{展望}：本文的模型框架具有较强的可移植性，未来可推广至以下方向：（1）其他呼吸道传染病（如COVID-19、RSV）；（2）其他人群密集场所（如中小学、企业、养老院）；（3）其他地区的高校（如北方城市、南方城市）；（4）多亚型竞争传播模型，考虑H1N1、H3N2、乙型流感的共循环效应．

\subsection{结论}

本文针对上海市高校甲型流感H3N2的传播特征与防控需求，建立了两群体分层SEIQR传播动力学模型，完成了参数估计、敏感性分析和防控策略优化．研究结果表明：（1）上海高校流感传播具有明显的冬夏双峰特征和开学效应，需采取差异化的季节性防控策略；（2）基础传播系数和学生-学生接触率是影响传播的关键因素，应优先干预；（3）当前疫苗接种率极低，需大幅提升至群体免疫阈值才能产生显著效果；（4）分场景、多目标的防控策略优化能够有效降低感染规模，同时平衡经济成本和教学影响．

本文的研究成果可为上海高校及疾控部门制定分季节、分场景的差异化防控策略提供定量决策支撑，助力构建健康、安全的校园环境．
